% --- Archon physics formalization source begin ---
% archon:physics
% archon:covers PhyXMiniProblems/problem_phyx_mini_0688.lean
% archon:source-report reports/phyx_mini/problem_phyx_mini_0688.source.json

\chapter{Physics problem phyx\_mini\_0688}
\label{ch:PhyXMiniProblems_problem_phyx_mini_0688}

\paragraph{Problem source.}
\#\# Physical scenario

Figure represents the total acceleration of a particle moving clockwise in a circle of radius \textbackslash{}( 2.50 \textbackslash{}text\{ m\} \textbackslash{}) at a certain instant of time.

\#\# Question

For that instant, find its tangential acceleration.

\#\# Answer choices

- A: A: \textbackslash{}( 2.50 \textbackslash{}text\{ m/s\}\textasciicircum{}2 \textbackslash{})

- B: B: \textbackslash{}( 10.00 \textbackslash{}text\{ m/s\}\textasciicircum{}2 \textbackslash{})

- C: C: \textbackslash{}( 7.50 \textbackslash{}text\{ m/s\}\textasciicircum{}2 \textbackslash{})

- D: D: \textbackslash{}( 5.00 \textbackslash{}text\{ m/s\}\textasciicircum{}2 \textbackslash{})

\#\# Figure caption (auxiliary; use the image as primary evidence)

The image illustrates a circular motion diagram. Here's a detailed description:

1. **Circle**: A dashed circle represents the path of the motion, indicating a circular trajectory.

2. **Radius**: A radius of 2.50 meters is shown, labeled on the circle with a line extending from the center to the edge of the circle.

3. **Object**: A small ball or point is located on the circle's perimeter where vectors are applied.

4. **Vectors**:
   - **Velocity (\textbackslash{}(\textbackslash{}vec\{v\}\textbackslash{}))**: Represented by a red arrow, tangent to the circle, pointing to the right.
   - **Acceleration (\textbackslash{}(\textbackslash{}vec\{a\}\textbackslash{}))**: Represented by a purple arrow, pointing inward at an angle.

5. **Angle**: An angle of 30.0° is marked between the radius and a line segment connected to the center of the circle.

6. **Acceleration Value**: The acceleration is labeled as \textbackslash{}(a = 15.0 \textbackslash{}, \textbackslash{}text\{m/s\}\textasciicircum{}2\textbackslash{}).

This diagram illustrates the relationship between velocity and acceleration vectors in circular motion. The angle indicates the direction of the vectors relative to the radius.

\#\# Dataset metadata

Rotational Motion; Spatial Relation Reasoning, Physical Model Grounding Reasoning

\paragraph{Recorded answer/context.}
C: C: \textbackslash{}( 7.50 \textbackslash{}text\{ m/s\}\textasciicircum{}2 \textbackslash{})

\paragraph{Figure/image path.}
/root/proposal\_for\_physic/science-mango/phyx\_mini\_run/phyx\_data/test\_image/688.png

\paragraph{Formalization target.}
create a compiling Lean file with sorry bodies at `PhyXMiniProblems/problem\_phyx\_mini\_0688.lean`. The Lean declarations must preserve the physical quantities, dimensions or dimensional roles, figure labels, governing-law hypotheses, and final relation expressed by this problem.
Use Mathlib/Physlib names found through LeanExplore where available. If a physics API is missing, introduce faithful local abstractions rather than scalar placeholder aliases.

\begin{theorem}[Physics formalization target]
\label{thm:physics:phyx_mini_0688:target}
\lean{PhyXMiniProblems.ProblemPhyXMini0688.tangentialAcceleration_is_7_point_50_metersPerSecondSquared}
\uses{def:physics:phyx-mini-0688:phyxminiproblems-problemphyxmini0688-accelerationinmeterspersecondsquared, def:physics:phyx-mini-0688:phyxminiproblems-problemphyxmini0688-circularmotionsnapshot, def:physics:phyx-mini-0688:phyxminiproblems-problemphyxmini0688-matchesclockwisecircularmotionscenario, def:physics:phyx-mini-0688:phyxminiproblems-problemphyxmini0688-matchessuppliedcircularaccelerationfigure, def:physics:phyx-mini-0688:phyxminiproblems-problemphyxmini0688-matchesprimaryfigurereadouts, def:physics:phyx-mini-0688:phyxminiproblems-problemphyxmini0688-hasphysicalcircularmotionparameters, def:physics:phyx-mini-0688:phyxminiproblems-problemphyxmini0688-satisfiescircularaccelerationdecomposition, def:physics:phyx-mini-0688:phyxminiproblems-problemphyxmini0688-matchesdisplayedtangentialacceleration, def:physics:phyx-mini-0688:phyxminiproblems-problemphyxmini0688-isuniquematchingtangentialaccelerationchoice}
The assigned autoformalize agent should translate this physics problem into Lean declarations in the covered file, with theorem and lemma proof bodies written as `by sorry`.
\end{theorem}
\begin{proof}
This is an autoformalization task, not a proof task. Produce statements that can later be proved by the physics prover without weakening the physical model.
\end{proof}

% --- Archon declaration topology begin ---
\section{Declaration topology}
\begin{definition}[speed Dimension]
  \label{def:physics:phyx-mini-0688:phyxminiproblems-problemphyxmini0688-speeddimension}
  \lean{PhyXMiniProblems.ProblemPhyXMini0688.speedDimension}
  The physical dimension of speed, L T⁻¹
\end{definition}
\begin{proof}
  This is the declared model definition of speed Dimension.
\end{proof}
\begin{definition}[acceleration Dimension]
  \label{def:physics:phyx-mini-0688:phyxminiproblems-problemphyxmini0688-accelerationdimension}
  \lean{PhyXMiniProblems.ProblemPhyXMini0688.accelerationDimension}
  The physical dimension of acceleration, L T⁻²
\end{definition}
\begin{proof}
  This is the declared model definition of acceleration Dimension.
\end{proof}
\begin{definition}[Length Magnitude]
  \label{def:physics:phyx-mini-0688:phyxminiproblems-problemphyxmini0688-lengthmagnitude}
  \lean{PhyXMiniProblems.ProblemPhyXMini0688.LengthMagnitude}
  A nonnegative, unit-independent physical length magnitude
\end{definition}
\begin{proof}
  This is the declared model definition of Length Magnitude.
\end{proof}
\begin{definition}[Speed Magnitude]
  \label{def:physics:phyx-mini-0688:phyxminiproblems-problemphyxmini0688-speedmagnitude}
  \lean{PhyXMiniProblems.ProblemPhyXMini0688.SpeedMagnitude}
  \uses{def:physics:phyx-mini-0688:phyxminiproblems-problemphyxmini0688-speeddimension}
  A nonnegative, unit-independent physical speed magnitude
\end{definition}
\begin{proof}
  This is the declared model definition of Speed Magnitude.
\end{proof}
\begin{definition}[Acceleration Magnitude]
  \label{def:physics:phyx-mini-0688:phyxminiproblems-problemphyxmini0688-accelerationmagnitude}
  \lean{PhyXMiniProblems.ProblemPhyXMini0688.AccelerationMagnitude}
  \uses{def:physics:phyx-mini-0688:phyxminiproblems-problemphyxmini0688-accelerationdimension}
  A nonnegative, unit-independent physical acceleration magnitude
\end{definition}
\begin{proof}
  This is the declared model definition of Acceleration Magnitude.
\end{proof}
\begin{definition}[length Magnitude Readout]
  \label{def:physics:phyx-mini-0688:phyxminiproblems-problemphyxmini0688-lengthmagnitudereadout}
  \lean{PhyXMiniProblems.ProblemPhyXMini0688.lengthMagnitudeReadout}
  \uses{def:physics:phyx-mini-0688:phyxminiproblems-problemphyxmini0688-lengthmagnitude}
  Read a physical length in a selected length unit
\end{definition}
\begin{proof}
  This is the declared model definition of length Magnitude Readout.
\end{proof}
\begin{definition}[speed Magnitude Readout]
  \label{def:physics:phyx-mini-0688:phyxminiproblems-problemphyxmini0688-speedmagnitudereadout}
  \lean{PhyXMiniProblems.ProblemPhyXMini0688.speedMagnitudeReadout}
  \uses{def:physics:phyx-mini-0688:phyxminiproblems-problemphyxmini0688-speedmagnitude}
  Read a speed magnitude in coherent selected length and time units
\end{definition}
\begin{proof}
  This is the declared model definition of speed Magnitude Readout.
\end{proof}
\begin{definition}[acceleration Magnitude Readout]
  \label{def:physics:phyx-mini-0688:phyxminiproblems-problemphyxmini0688-accelerationmagnitudereadout}
  \lean{PhyXMiniProblems.ProblemPhyXMini0688.accelerationMagnitudeReadout}
  \uses{def:physics:phyx-mini-0688:phyxminiproblems-problemphyxmini0688-accelerationmagnitude}
  Read an acceleration magnitude in coherent selected units
\end{definition}
\begin{proof}
  This is the declared model definition of acceleration Magnitude Readout.
\end{proof}
\begin{definition}[length In Meters]
  \label{def:physics:phyx-mini-0688:phyxminiproblems-problemphyxmini0688-lengthinmeters}
  \lean{PhyXMiniProblems.ProblemPhyXMini0688.lengthInMeters}
  \uses{def:physics:phyx-mini-0688:phyxminiproblems-problemphyxmini0688-lengthmagnitude, def:physics:phyx-mini-0688:phyxminiproblems-problemphyxmini0688-lengthmagnitudereadout}
  Metre readout of a physical length
\end{definition}
\begin{proof}
  This is the declared model definition of length In Meters.
\end{proof}
\begin{definition}[speed In Meters Per Second]
  \label{def:physics:phyx-mini-0688:phyxminiproblems-problemphyxmini0688-speedinmeterspersecond}
  \lean{PhyXMiniProblems.ProblemPhyXMini0688.speedInMetersPerSecond}
  \uses{def:physics:phyx-mini-0688:phyxminiproblems-problemphyxmini0688-speedmagnitude, def:physics:phyx-mini-0688:phyxminiproblems-problemphyxmini0688-speedmagnitudereadout}
  Metres-per-second readout of a speed magnitude
\end{definition}
\begin{proof}
  This is the declared model definition of speed In Meters Per Second.
\end{proof}
\begin{definition}[acceleration In Meters Per Second Squared]
  \label{def:physics:phyx-mini-0688:phyxminiproblems-problemphyxmini0688-accelerationinmeterspersecondsquared}
  \lean{PhyXMiniProblems.ProblemPhyXMini0688.accelerationInMetersPerSecondSquared}
  \uses{def:physics:phyx-mini-0688:phyxminiproblems-problemphyxmini0688-accelerationmagnitude, def:physics:phyx-mini-0688:phyxminiproblems-problemphyxmini0688-accelerationmagnitudereadout}
  Metres-per-second-squared readout of an acceleration magnitude
\end{definition}
\begin{proof}
  This is the declared model definition of acceleration In Meters Per Second Squared.
\end{proof}
\begin{definition}[Rotation Sense]
  \label{def:physics:phyx-mini-0688:phyxminiproblems-problemphyxmini0688-rotationsense}
  \lean{PhyXMiniProblems.ProblemPhyXMini0688.RotationSense}
  Sense of motion around the circular path
\end{definition}
\begin{proof}
  This is the declared model definition of Rotation Sense.
\end{proof}
\begin{definition}[Particle Ray]
  \label{def:physics:phyx-mini-0688:phyxminiproblems-problemphyxmini0688-particleray}
  \lean{PhyXMiniProblems.ProblemPhyXMini0688.ParticleRay}
  Named rays based at the pictured particle. The total-acceleration ray lies between the inward radial ray and the clockwise tangential ray in the image
\end{definition}
\begin{proof}
  This is the declared model definition of Particle Ray.
\end{proof}
\begin{definition}[Figure Vector]
  \label{def:physics:phyx-mini-0688:phyxminiproblems-problemphyxmini0688-figurevector}
  \lean{PhyXMiniProblems.ProblemPhyXMini0688.FigureVector}
  The two vector arrows explicitly drawn in the supplied image
\end{definition}
\begin{proof}
  This is the declared model definition of Figure Vector.
\end{proof}
\begin{definition}[Figure Label]
  \label{def:physics:phyx-mini-0688:phyxminiproblems-problemphyxmini0688-figurelabel}
  \lean{PhyXMiniProblems.ProblemPhyXMini0688.FigureLabel}
  Literal or mathematical labels visible in the supplied image
\end{definition}
\begin{proof}
  This is the declared model definition of Figure Label.
\end{proof}
\begin{definition}[Figure Feature]
  \label{def:physics:phyx-mini-0688:phyxminiproblems-problemphyxmini0688-figurefeature}
  \lean{PhyXMiniProblems.ProblemPhyXMini0688.FigureFeature}
  Geometric features visible in the circular-motion diagram
\end{definition}
\begin{proof}
  This is the declared model definition of Figure Feature.
\end{proof}
\begin{definition}[Velocity Vector At Particle]
  \label{def:physics:phyx-mini-0688:phyxminiproblems-problemphyxmini0688-velocityvectoratparticle}
  \lean{PhyXMiniProblems.ProblemPhyXMini0688.VelocityVectorAtParticle}
  \uses{def:physics:phyx-mini-0688:phyxminiproblems-problemphyxmini0688-speedmagnitude, def:physics:phyx-mini-0688:phyxminiproblems-problemphyxmini0688-particleray}
  A physical velocity vector at the instant, represented by its dimensionful magnitude and its geometrically named direction ray
\end{definition}
\begin{proof}
  This is the declared model definition of Velocity Vector At Particle.
\end{proof}
\begin{definition}[Acceleration Vector At Particle]
  \label{def:physics:phyx-mini-0688:phyxminiproblems-problemphyxmini0688-accelerationvectoratparticle}
  \lean{PhyXMiniProblems.ProblemPhyXMini0688.AccelerationVectorAtParticle}
  \uses{def:physics:phyx-mini-0688:phyxminiproblems-problemphyxmini0688-accelerationmagnitude, def:physics:phyx-mini-0688:phyxminiproblems-problemphyxmini0688-particleray}
  A physical acceleration vector at the instant, represented by its dimensionful magnitude and its geometrically named direction ray
\end{definition}
\begin{proof}
  This is the declared model definition of Acceleration Vector At Particle.
\end{proof}
\begin{definition}[Circular Acceleration Figure]
  \label{def:physics:phyx-mini-0688:phyxminiproblems-problemphyxmini0688-circularaccelerationfigure}
  \lean{PhyXMiniProblems.ProblemPhyXMini0688.CircularAccelerationFigure}
  \uses{def:physics:phyx-mini-0688:phyxminiproblems-problemphyxmini0688-particleray, def:physics:phyx-mini-0688:phyxminiproblems-problemphyxmini0688-figurevector, def:physics:phyx-mini-0688:phyxminiproblems-problemphyxmini0688-figurelabel, def:physics:phyx-mini-0688:phyxminiproblems-problemphyxmini0688-figurefeature}
  Primary-raster information from image 688.png. Numerical marker fields are scalar figure readouts; premises below relate them to independent physical quantities in the setup
\end{definition}
\begin{proof}
  This is the declared model definition of Circular Acceleration Figure.
\end{proof}
\begin{definition}[Circular Motion Snapshot]
  \label{def:physics:phyx-mini-0688:phyxminiproblems-problemphyxmini0688-circularmotionsnapshot}
  \lean{PhyXMiniProblems.ProblemPhyXMini0688.CircularMotionSnapshot}
  \uses{def:physics:phyx-mini-0688:phyxminiproblems-problemphyxmini0688-lengthmagnitude, def:physics:phyx-mini-0688:phyxminiproblems-problemphyxmini0688-accelerationmagnitude, def:physics:phyx-mini-0688:phyxminiproblems-problemphyxmini0688-rotationsense, def:physics:phyx-mini-0688:phyxminiproblems-problemphyxmini0688-velocityvectoratparticle, def:physics:phyx-mini-0688:phyxminiproblems-problemphyxmini0688-accelerationvectoratparticle, def:physics:phyx-mini-0688:phyxminiproblems-problemphyxmini0688-circularaccelerationfigure}
  Physical state of the particle at the depicted instant. The two component magnitudes are observables constrained by the kinematic laws below; neither is a local expansion of an answer choice
\end{definition}
\begin{proof}
  This is the declared model definition of Circular Motion Snapshot.
\end{proof}
\begin{definition}[Matches Clockwise Circular Motion Scenario]
  \label{def:physics:phyx-mini-0688:phyxminiproblems-problemphyxmini0688-matchesclockwisecircularmotionscenario}
  \lean{PhyXMiniProblems.ProblemPhyXMini0688.MatchesClockwiseCircularMotionScenario}
  \uses{def:physics:phyx-mini-0688:phyxminiproblems-problemphyxmini0688-circularmotionsnapshot}
  The clockwise circular-motion facts stated in the problem and diagram
\end{definition}
\begin{proof}
  This is the declared model definition of Matches Clockwise Circular Motion Scenario.
\end{proof}
\begin{definition}[Matches Supplied Circular Acceleration Figure]
  \label{def:physics:phyx-mini-0688:phyxminiproblems-problemphyxmini0688-matchessuppliedcircularaccelerationfigure}
  \lean{PhyXMiniProblems.ProblemPhyXMini0688.MatchesSuppliedCircularAccelerationFigure}
  \uses{def:physics:phyx-mini-0688:phyxminiproblems-problemphyxmini0688-circularmotionsnapshot}
  Incidence and spatial information read directly from the supplied image. In particular, the marked angle runs from the inward radius to the total acceleration arrow, not from the tangent to that arrow
\end{definition}
\begin{proof}
  This is the declared model definition of Matches Supplied Circular Acceleration Figure.
\end{proof}
\begin{definition}[Matches Primary Figure Readouts]
  \label{def:physics:phyx-mini-0688:phyxminiproblems-problemphyxmini0688-matchesprimaryfigurereadouts}
  \lean{PhyXMiniProblems.ProblemPhyXMini0688.MatchesPrimaryFigureReadouts}
  \uses{def:physics:phyx-mini-0688:phyxminiproblems-problemphyxmini0688-lengthinmeters, def:physics:phyx-mini-0688:phyxminiproblems-problemphyxmini0688-accelerationinmeterspersecondsquared, def:physics:phyx-mini-0688:phyxminiproblems-problemphyxmini0688-circularmotionsnapshot}
  The three numerical readouts shown in the primary image, together with their interpretation as physical radius, total-acceleration magnitude, and radian angle. No tangential-acceleration value appears here
\end{definition}
\begin{proof}
  This is the declared model definition of Matches Primary Figure Readouts.
\end{proof}
\begin{definition}[Has Physical Circular Motion Parameters]
  \label{def:physics:phyx-mini-0688:phyxminiproblems-problemphyxmini0688-hasphysicalcircularmotionparameters}
  \lean{PhyXMiniProblems.ProblemPhyXMini0688.HasPhysicalCircularMotionParameters}
  \uses{def:physics:phyx-mini-0688:phyxminiproblems-problemphyxmini0688-lengthinmeters, def:physics:phyx-mini-0688:phyxminiproblems-problemphyxmini0688-speedinmeterspersecond, def:physics:phyx-mini-0688:phyxminiproblems-problemphyxmini0688-accelerationinmeterspersecondsquared, def:physics:phyx-mini-0688:phyxminiproblems-problemphyxmini0688-circularmotionsnapshot}
  Positivity and acute-angle conditions selecting the pictured branch
\end{definition}
\begin{proof}
  This is the declared model definition of Has Physical Circular Motion Parameters.
\end{proof}
\begin{definition}[Satisfies Circular Acceleration Decomposition]
  \label{def:physics:phyx-mini-0688:phyxminiproblems-problemphyxmini0688-satisfiescircularaccelerationdecomposition}
  \lean{PhyXMiniProblems.ProblemPhyXMini0688.SatisfiesCircularAccelerationDecomposition}
  \uses{def:physics:phyx-mini-0688:phyxminiproblems-problemphyxmini0688-lengthmagnitudereadout, def:physics:phyx-mini-0688:phyxminiproblems-problemphyxmini0688-speedmagnitudereadout, def:physics:phyx-mini-0688:phyxminiproblems-problemphyxmini0688-accelerationmagnitudereadout, def:physics:phyx-mini-0688:phyxminiproblems-problemphyxmini0688-circularmotionsnapshot}
  Instantaneous circular-motion kinematics. The radial component is the centripetal acceleration v²/r. Resolving the total acceleration at angle phi from the inward radius gives a\_r = a cos(phi) and a\_t = a sin(phi). These are general laws and contain no numerical answer or answer label
\end{definition}
\begin{proof}
  This is the declared model definition of Satisfies Circular Acceleration Decomposition.
\end{proof}
\begin{definition}[Answer Choice]
  \label{def:physics:phyx-mini-0688:phyxminiproblems-problemphyxmini0688-answerchoice}
  \lean{PhyXMiniProblems.ProblemPhyXMini0688.AnswerChoice}
  Labels of the four tangential-acceleration choices in the source
\end{definition}
\begin{proof}
  This is the declared model definition of Answer Choice.
\end{proof}
\begin{definition}[displayed Tangential Acceleration In Meters Per Second Squared]
  \label{def:physics:phyx-mini-0688:phyxminiproblems-problemphyxmini0688-displayedtangentialaccelerationinmeterspersecondsquared}
  \lean{PhyXMiniProblems.ProblemPhyXMini0688.displayedTangentialAccelerationInMetersPerSecondSquared}
  \uses{def:physics:phyx-mini-0688:phyxminiproblems-problemphyxmini0688-answerchoice}
  Metres-per-second-squared value printed beside each answer label
\end{definition}
\begin{proof}
  This is the declared model definition of displayed Tangential Acceleration In Meters Per Second Squared.
\end{proof}
\begin{definition}[recorded Dataset Answer]
  \label{def:physics:phyx-mini-0688:phyxminiproblems-problemphyxmini0688-recordeddatasetanswer}
  \lean{PhyXMiniProblems.ProblemPhyXMini0688.recordedDatasetAnswer}
  \uses{def:physics:phyx-mini-0688:phyxminiproblems-problemphyxmini0688-answerchoice}
  The answer label recorded by the source dataset, retained as metadata
\end{definition}
\begin{proof}
  This is the declared model definition of recorded Dataset Answer.
\end{proof}
\begin{definition}[Matches Displayed Tangential Acceleration]
  \label{def:physics:phyx-mini-0688:phyxminiproblems-problemphyxmini0688-matchesdisplayedtangentialacceleration}
  \lean{PhyXMiniProblems.ProblemPhyXMini0688.MatchesDisplayedTangentialAcceleration}
  \uses{def:physics:phyx-mini-0688:phyxminiproblems-problemphyxmini0688-accelerationinmeterspersecondsquared, def:physics:phyx-mini-0688:phyxminiproblems-problemphyxmini0688-circularmotionsnapshot, def:physics:phyx-mini-0688:phyxminiproblems-problemphyxmini0688-answerchoice, def:physics:phyx-mini-0688:phyxminiproblems-problemphyxmini0688-displayedtangentialaccelerationinmeterspersecondsquared}
  A choice agrees exactly with the physical tangential-acceleration readout
\end{definition}
\begin{proof}
  This is the declared model definition of Matches Displayed Tangential Acceleration.
\end{proof}
\begin{definition}[Is Unique Matching Tangential Acceleration Choice]
  \label{def:physics:phyx-mini-0688:phyxminiproblems-problemphyxmini0688-isuniquematchingtangentialaccelerationchoice}
  \lean{PhyXMiniProblems.ProblemPhyXMini0688.IsUniqueMatchingTangentialAccelerationChoice}
  \uses{def:physics:phyx-mini-0688:phyxminiproblems-problemphyxmini0688-circularmotionsnapshot, def:physics:phyx-mini-0688:phyxminiproblems-problemphyxmini0688-answerchoice, def:physics:phyx-mini-0688:phyxminiproblems-problemphyxmini0688-matchesdisplayedtangentialacceleration}
  A choice is the unique displayed value matching the physical result
\end{definition}
\begin{proof}
  This is the declared model definition of Is Unique Matching Tangential Acceleration Choice.
\end{proof}
% --- Archon declaration topology end ---
% --- Archon physics formalization source end ---
